\documentclass[11pt,a4paper]{article}
\usepackage[utf8]{inputenc}
\usepackage{graphicx}
\usepackage[left=2.5cm,top=3cm,right=2.5cm,bottom=3cm,bindingoffset=0.5cm]{geometry}
\usepackage{AEDLogica, AEDEspecificacion, AEDTADs}
\usepackage{caratula}
\newlength{\miSangria}
\setlength{\miSangria}{2em}

% Asi pueden escribir nuevos comandos. 
% Este por ejemplo asegura q los nombres 
% que figuren con una tipografia diferenciada  
\newcommand{\Tipo}[1]{\mathsf{#1}}
% la sintaxis es \newcommand{\nombreDeLaMacro}[cantidadDeParametros]{Lo que va ser remplazado por el macro} 
\newcommand{\norm}[1]{\vert#1\vert}


\begin{document}

\section{Consultas y correcciones internas}
\vspace{1.0cm}

\textbf{AcumularMillas}
\begin{enumerate}
    \item AcumularMillas:
    \begin{enumerate}
        \item Consultar que se hace cuando las millas restantes de la promocion son menores que la distancia. (La decision tomada fue no hacer nada, consultar resolucion).
        \item Consultar los entoncesluego y los yluego
        \item Consultar para las funciones EsVuelo, EsTransferencia, EsTransferenciaEnviada, principalmente si tengo que verificar todos los campos de antes o no (actualizacion creo que no porque es un pred?)
        \item Consultar anidacion de ifthenelse que en el grupo hablaron que era medio dudoso. Para mi era valido
    \end{enumerate}
    \item CanjearMillas.
    \begin{enumerate}
        \item En la practica nunca usamos una metavarialbe asignada en el requiere para hacer esto de millasDisponibles = algo, para eso esta la funcion millasDisponibles que es un pred que esta creado en la funcion acumular millas.
    \end{enumerate}
    \item Tercer elemento.
\end{enumerate}

\end{document}